\section{Related Work}

\subsection{Spurious Correlations and Group Robustness}

Spurious correlations cause models to rely on dataset biases rather than causal features. \citet{sagawa2020distributionally} introduced Group DRO, which minimizes worst-case loss over predefined groups, achieving strong results on Waterbirds and CelebA. However, Group DRO requires group annotations during training --- expensive and often unavailable.

Methods without group annotations have emerged. \textbf{Just Train Twice (JTT)} \citep{liu2021just} trains an initial ERM model, identifies misclassified examples (which correlate with minority groups), and upweights them in a second training run. \textbf{Learning from Failure (LfF)} \citep{nam2020learning} explicitly trains a biased network and uses its failures to guide debiasing. Both methods exploit the insight that spurious features are learned early, but require two-stage training.

\textbf{Deep Feature Reweighting (DFR)} \citep{kirichenko2022last} demonstrates that pretrained ERM features are sufficient for state-of-the-art worst-group accuracy when the last layer is retrained on group-balanced data. This finding directly informs our negative result: if CLIP features already separate spurious and core concepts, there are no emergence dynamics left to observe.

\subsection{Simplicity Bias and Learning Dynamics}

\citet{shah2020pitfalls} established that SGD exhibits \textbf{simplicity bias}: networks learn the simplest predictive features first and may never learn more complex features even when they have higher predictive power. This explains why spurious (often simpler) correlations dominate.

\textbf{Gradient Starvation} \citep{pezeshki2021gradient} provides a dynamical systems perspective: cross-entropy minimization on features with different frequencies causes some features to receive diminishing gradient signal. Once simple features achieve low loss, complex features are ``starved'' of learning signal.

Our work extends this literature by testing whether emergence \emph{uniformity} (variance across sample subsets) differs between spurious and core features. The negative result suggests that while emergence \emph{timing} may differ during training, this signal is lost in pretrained representations.

\subsection{Linear Probing for Representation Analysis}

Linear probes are standard tools for analyzing pretrained representations \citep{alain2017understanding}. The CLIP evaluation protocol \citep{radford2021learning} uses logistic regression with regularization-strength sweeps (C-sweep) to assess representation quality.

We adapted this protocol for emergence dynamics: treating C as an epoch proxy and computing CV across sample subsets. Our failure reveals a category error: C-sweep produces representation quality scores, not learning trajectory data. Convex optimization converges in a single pass regardless of C; there is no ``emergence'' to measure.

\subsection{Positioning Our Work}

\begin{table}[h]
\centering
\begin{tabular}{@{}lccl@{}}
\toprule
Method & Stage & Annotations & Detects During Training \\
\midrule
Group DRO & 1 & Yes & N/A \\
JTT & 2 & No & Yes (early errors) \\
LfF & 2 & No & Yes (biased network) \\
DFR & 1 & Yes (balanced retrain) & No \\
\textbf{EUR (proposed)} & 1 & No & \textbf{Intended: Yes} \\
\bottomrule
\end{tabular}
\end{table}

Our contribution is a negative result: the CV-based detection that EUR depends on does not work on frozen pretrained features.
